% !TEX program = lualatex
\documentclass[a4paper,11pt]{scrartcl}

\usepackage[
  lang=zh,
  mode=journal,
  theme=graphite,
  toc=true,
  toc-layout=page
]{synaptic}
\usepackage{booktabs,tabularx}
\newcolumntype{Y}{>{\raggedright\arraybackslash}X}
\newcommand{\DocCmd}[1]{\textcolor{syn-primary!80!black}{\texttt{\textbackslash#1}}}

\SynapticHeader{synaptic 技术参考}
\SynapticSubHeader{架构、不变量与发布流程}
\SynapticPreprint{v5.3.0}
\SynapticTitle{synaptic 技术参考}
\SynapticAuthor{synaptic 项目}
\SynapticAffiliation{开源 LaTeX 社区}
\SynapticSubject{synaptic 5.3.0 技术参考}
\SynapticKeywords{LaTeX3, 架构, l3build, 回归测试}

\begin{synabstract}
  本文档说明 synaptic 的模块图、加载期不变量、命名空间规则、字体与颜色子系统、
  扩展点、生成文件策略以及自动化验证流程。
\end{synabstract}

\begin{document}
\SynapticMakeTitle

\section{架构}

\begin{verbatim}
synaptic.sty
  `-- synaptic-base.sty
      |-- synaptic-lang-{en,zh}.def
      |-- synaptic-color.sty
      |-- synaptic-fonts.sty
      |-- synaptic-layout.sty
      |-- synaptic-theorem.sty
      |-- synaptic-title.sty       公共元数据与渲染器
      `-- 模式模块
          |-- synaptic-journal.sty     journal
          |-- synaptic-book.sty        book
          |-- synaptic-lecture.sty     lecture
          |-- synaptic-notes.sty       notes
          `-- synaptic-beamer.sty      beamer
\end{verbatim}

薄封装层把用户选项转发给 \texttt{synaptic-base}，并提供带命名空间的致谢命令。
基础模块负责检查引擎与文档类、解析选项、加载标签、验证结构要求，然后分发核心
模块和模式模块。语言定义文件先于消费这些标签的模块加载。

\subsection{模块职责}

\noindent{\setlength{\parindent}{0pt}\noindent\begin{tabularx}{\textwidth}{@{}l l Y@{}}
  \toprule
  \textbf{模块} & \textbf{类别} & \textbf{职责} \\
  \midrule
  \texttt{synaptic.sty} & 封装 & 选项转发、致谢命令 \\
  \texttt{synaptic-base.sty} & 核心 & 引擎/文档类检查、键系统、分发 \\
  \texttt{synaptic-lang-*.def} & 数据 & 双语标签常量 \\
  \texttt{synaptic-color.sty} & 核心 & 主题登记、语义令牌、交付模式 \\
  \texttt{synaptic-fonts.sty} & 核心 & 固定的西文与中文字体栈 \\
  \texttt{synaptic-layout.sty} & 核心 & 页面几何、排版、标题、编号 \\
  \texttt{synaptic-theorem.sty} & 核心 & 定理样式、环境、证明 \\
  \texttt{synaptic-title.sty} & 核心 & 元数据模型与五种标题渲染器 \\
  \texttt{synaptic-journal.sty} & 模式 & 期刊导航与声明 \\
  \texttt{synaptic-book.sty} & 模式 & 章/部分样式、前后文命令 \\
  \texttt{synaptic-lecture.sty} & 模式 & 信息条、教学卡片、导航 \\
  \texttt{synaptic-notes.sty} & 模式 & 笔记信息条与知识卡 \\
  \texttt{synaptic-beamer.sty} & 模式 & 幻灯片模板、进度、目录/过渡帧 \\
  \bottomrule
\end{tabularx}}

\noindent{\setlength{\parindent}{0pt}\noindent\begin{tabularx}{\textwidth}{@{}l l Y@{}}
  \toprule
  \textbf{模式} & \textbf{文档类不变量} & \textbf{附加模块} \\
  \midrule
  journal & KOMA 文章/报告类 & journal \\
  book & 含章和前文接口的 KOMA 类 & book \\
  lecture & KOMA 文章/报告类 & lecture \\
  notes & KOMA 文章/报告类 & notes \\
  beamer & Beamer 文档类 & beamer \\
  \bottomrule
\end{tabularx}}

\section{配置生命周期}

宏包级 \texttt{synaptic} 键系统在模块加载前执行。模式、语言、标题布局、
版心尺度、密度、编号、输出颜色模式、定理样式和装订偏移因此成为结构不变量。
运行时键系统只暴露无需重载宏包就能应用的
\texttt{theme}、\texttt{background}、\texttt{toc} 和 \texttt{toc-layout}。运行时请求中的结构键与未知键
直接报错；宏包选项中的未知键同样报错。

书籍模式要求章和前文接口同时存在。文档类契约在选项解析后检查：Beamer 模式要求
Beamer 文档类，四种纸面模式要求 KOMA-Script。LuaLaTeX 与文档类要求在不满足时立即给出
致命诊断，而不是延迟到后续的未定义命令。

\section{设计不变量}

以下性质由实现有意保证，并由回归套件守护。

\begin{itemize}
  \item \textbf{确定性排版}。固定字体栈，行距、版心与标题尺寸全部显式，
    不因机器上的额外字体而改变设计。
  \item \textbf{不全局屏蔽坏盒子}。emergencystretch 与 tolerance 经过调校，
    溢出与欠满警告保留给作者。
  \item \textbf{不强制纸张、不重置页码}。几何布局跟随文档类；期刊标题不会
    悄悄重新开始页码。
  \item \textbf{只做一次 Unicode 转换}。PDF 元数据以源词元交给 hyperref；
    预先转换会造成双重编码。
  \item \textbf{公开名所有权明确}。所有公开命令与环境都带命名空间；
    文档类或其他宏包提供的通用环境不会被修改。
  \item \textbf{语义颜色架构}。组件只消费语义令牌；主题与交付模式只改变
    令牌，不改变组件逻辑。
\end{itemize}

\section{子系统}

\subsection{颜色令牌}

内置主题登记在全局属性表中。启用主题时，主色、辅色与强调色先复制到稳定基色角色，
再派生公开显示色、三档表面、边框、分隔线、链接、可读弱化文字、装饰浅色与正文色。页面背景键
\texttt{background}（任意 xcolor 表达式）成为 \texttt{syn-bg} 基础令牌；所有
派生角色与它混合，并按相对亮度切换：背景偏暗时中性文字翻转成浅色墨，品牌色、强调色与状态色同步提亮。
稳定基色不被覆写，因而后续背景切换可逆。print 与 mono 输出模式只改变同一派生过程，不改变组件代码。
危险、警告和成功色保留独立于品牌配色的稳定语义基色。

\begin{verbatim}
\SynapticDefineTheme{名称}{主色}{辅色}[强调色]
\SynapticSetup{theme=名称}
\end{verbatim}

颜色参数是 xcolor 表达式。主题名必须是经过校验的小写标识符；定义不可变，
启用操作显式完成。

\subsection{字体选择}

字体探测只负责校验固定字体栈：西文正文与数学使用 Libertinus Serif、Sans、Mono
和 Math；其粗体数学版本由同一 Libertinus Math 字面确定性合成。\texttt{lang=zh}
以无预设字体方案加载 ctex，中文使用思源宋体/黑体 CN，
中文强调使用霞鹜文楷。缺少字体直接报错，不触发另一套视觉设计。

\subsection{版式与文档结构}

页面布局结合模式默认值、版心尺度与密度档位。journal 与 lecture 优先适中行长，notes
支持紧凑单栏/双栏，book 使用镜像内外边、装订补偿、外侧页眉和空白偶数页。孤行惩罚、
垂直节奏、列表间距、行距与段首缩进都是明确配置，
不再用全局方式隐藏坏盒子。

Beamer 走独立分支：由文档类保持对幻灯片画布的所有权，\texttt{measure} 映射为对称帧边距，
\texttt{density} 映射为演示行距。KOMA 专用的几何、图表标题、列表、结构标题与目录声明不会在
此分支执行。模式模块安装语义 Beamer 颜色、可变高信号导轨帧标题、对齐的带编号目录项、单层原生卡片、列表标记、
带轨道且安全钳制的进度页脚，以及显式的目录与章节过渡帧命令。定理模块提供一个共享卡片原语，由标题元数据与各模式卡片继承。

KOMA 标题字体与间距在 \texttt{begindocument/before} 钩子中安装，所有编号层级共用
无衬线标题家族：\texttt{section} 用主色，\texttt{subsection} 用深色，
\texttt{subsubsection} 更小且略微淡化，使层级一目了然。hyperref、bookmark、
microtype、图表标题标点、列表与模式化公式编号由布局模块统一配置。各模式模块通过共享
辅助函数安装页眉发丝线（\texttt{headsepline}），并用专用的 KOMA
\texttt{headsepline} 字体着色，使装饰线比页眉文字更浅。

目录通过 KOMA 的 \texttt{\string\DeclareTOCStyleEntry} 在其输出前立即细化：顶层条目
以加粗主色重现将页面标题颜色，深层条目保持深色，并收紧编号对齐与垂直节奏。图表标题
标签使用次要主题色。

\subsection{标题与元数据}

标题值是贯穿文档的状态，因此保存在 expl3 全局词元表和序列中。作者与机构标记使用
平行序列，机构与邮箱可重复。
公共元数据分别驱动期刊内联标题、书籍扉页序列、课程信息条、紧凑笔记信息条与 Beamer 标题帧。标题命令
验证非空标题并控制可选目录位置，不再隐式重置用户页码。带命名空间的
\texttt{synabstract} 捕获摘要；标准 \texttt{abstract} 与 \verb|\maketitle| 均不修改。

PDF 信息以源词元传给 hyperref，由 hyperref 只执行一次 Unicode PDF 字符串转换。预先
\verb|\pdfstringdef| 再传给 \verb|\hypersetup| 会导致双重编码，因此已禁止这种路径。
Beamer 模式中，带命名空间的设置命令会在 \verb|\begin{document}| 前同步到文档类原生的
标题、作者、主题和关键词接口。由于 Beamer 在文档开始时冻结这些字段，共享的后期元数据写入器会
在此模式中跳过。

\subsection{陈述}

thmtools 在 amsthm 之上定义 plain、definition 和 remark 三种样式。内置陈述族
（\texttt{syntheorem}、\texttt{synlemma}、\texttt{synproposition}、
\texttt{syndefinition}、\texttt{syncorollary}、\texttt{synaxiom}、\texttt{synconjecture}、
\texttt{synclaim}、\texttt{synproperty}、\texttt{synexample}、\texttt{synexercise}、
\texttt{syncriterion}、\texttt{synconvention}、\texttt{synproblem}、\texttt{synsolution}、
\texttt{synremark}、\texttt{synnote} 与 \texttt{synwarning}）使用模式化共享计数器：
书籍默认按章，其他模式默认按节。在 \texttt{theorem-style=boxed} 之下，编号的陈述
环境在 \texttt{\string\begin{document}} 处由 tcolorbox 包装（主题色背景配西侧强调条）；
\texttt{synproof} 保持无框，且不会重定义标准 \texttt{proof}。
\texttt{theorem-style=plain} 关闭陈述包装。例题与练习为
plain 样式陈述，警告使用注记呈现。由于所有陈述共享定理计数器，本模块还会按真实结构
父级重建每个环境的超链接锚点（\texttt{theH<名称>}），保证跨章跨节的链接唯一。

\section{命名空间规则}

\noindent{\setlength{\parindent}{0pt}\noindent\begin{tabularx}{\textwidth}{@{}l l Y@{}}
  \toprule
  \textbf{范围} & \textbf{模式} & \textbf{示例} \\
  \midrule
  公开命令 & \texttt{\textbackslash Synaptic\ldots} & \DocCmd{SynapticDefineTheme} \\
  公开环境 & \texttt{syn...} & \texttt{synsummary} \\
  公开颜色令牌 & \texttt{syn-...} & \texttt{syn-primary} \\
  内部函数 & \texttt{\textbackslash\_\_synaptic\_\ldots:} & 标题渲染辅助函数 \\
  私有局部状态 & \texttt{\textbackslash l\_\_synaptic\_\ldots} & 标题布局临时值 \\
  私有全局状态 & \texttt{\textbackslash g\_\_synaptic\_\ldots} & 模式、元数据或注册表 \\
  私有常量 & \texttt{\textbackslash c\_\_synaptic\_\ldots} & 标签或环境集合 \\
  \bottomrule
\end{tabularx}}

expl3 的 \texttt{g} 与 \texttt{l} 前缀只描述生命周期，不描述所有权；所有权由
\texttt{synaptic} 模块名与私有双下划线独立表达。因此文档级状态仍使用私有的
\texttt{g\_\_synaptic\_...}，短期临时值使用 \texttt{l\_\_synaptic\_...}。

新公开名必须遵循同样的所有权规则。\verb|\SynapticNewTheorem| 接受的自定义
环境名以 \texttt{syn} 开头且只含小写字母与数字；未命名空间化或重名都会报错。

\section{扩展点}

\begin{itemize}
  \item 在新的 \texttt{synaptic-lang-*.def} 中增加标签，再扩展语言键与定义文件分发。
  \item 新字体组需同时增加完整性谓词、设置函数、显式分支，并明确其在自动链中的位置。
  \item 新模式需扩展基础键系统、文档类验证、布局分发、标题/目录分发和模块加载器，
    并声明、测试必需的文档类能力。
  \item 视觉扩展应基于语义颜色令牌，而不是固定 RGB 值。
  \item 每项扩展都应同步增加回归测试；测试套件是上述不变量的守护者。
\end{itemize}

\section{自动化验证}

\texttt{synaptic.dtx} 是唯一源头。DocStrip 通过 \texttt{synaptic.ins} 生成可安装的
\texttt{.sty} 和 \texttt{.def} 文件到 \texttt{build/unpacked}。提取模块、编译 PDF
与发布压缩包均不跟踪，因此实现只有一个需要审阅和更新的事实来源。

回归测试的组织方式如下：

{\setlength{\parindent}{0pt}\noindent\begin{tabularx}{\textwidth}{@{}l Y@{}}
  \toprule
  \textbf{测试组} & \textbf{覆盖内容} \\
  \midrule
  \texttt{journal}、\texttt{lecture}、\texttt{notes}、\texttt{book}、\texttt{beamer} & 核心模式
  行为、导航与默认值 \\
  \texttt{book-zh}、\texttt{lecture-zh}、\texttt{journal-zh}、\texttt{beamer-zh}、
  \texttt{language-zh} & 中文标签、CJK 字体与双语层次 \\
  \texttt{themes} & 五种内置主题与自定义主题循环 \\
  \texttt{color-mode-print}、\texttt{options} & 打印与单色交付模式 \\
  \texttt{customization} & 运行时设置与自定义定理 \\
  \texttt{toc-layouts}、\texttt{title-layouts} & 目录与标题构成 \\
  \texttt{layout-profiles} & 密度、版心尺度、装订偏移 \\
  \texttt{theorem-boxes}、\texttt{strict-options} & 完整陈述族、命名空间验证与严格选项错误 \\
  \texttt{book-full} & 完整标题序列与前后文流程 \\
  \texttt{lecture-continuous} & 连续编号档案 \\
  \bottomrule
\end{tabularx}}

每组使用标准引擎排版一份文档，并将规范化日志与已提交的期望文件比较。共享入口
\texttt{scripts/verify} 运行套件、构建源码文档、四份手册和所有示例，并在出现警告、
溢出盒子、缺字符或 TeX 错误时失败；CI 调用同一入口。

\section{源码与发布流程}

\begin{verbatim}
./scripts/verify  # 测试、手册、示例与诊断
l3build unpack    # 提取 .sty/.def 到 build/unpacked
l3build ctan      # 生成经过测试的 TDS 与 CTAN 压缩包
l3build install   # 复制宏包文件到用户 texmf 目录树
l3build uninstall # 移除先前安装
\end{verbatim}

\texttt{ctan} 目标生成 \texttt{synaptic-ctan.zip}（源码树与生成文档、示例）以及内嵌的
\texttt{synaptic.tds.zip}（TDS 布局）。安装会就地替换旧版本。仓库不保留兼容副本
或迁移垫片；已移除接口统一记录在变更日志中。

\end{document}
